% Printed source list; use the main edition preamble. No invented source column headings.
% AB01-PDF0072-body-L033
\SourceLine{AB01-PDF0072-body-L033}{’ ({\itshape hamzah}) in medio et fine vocabulorum asperam vocis interruptionem designat.}
% AB01-PDF0072-body-L034
\SourceLine{AB01-PDF0072-body-L034}{‘ (\textarabic{\upshape ع}) est gutturalis surda, nostris linguis ignota.}
% AB01-PDF0072-body-L035
\SourceLine{AB01-PDF0072-body-L035}{{\itshape dh} (\textarabic{\upshape ذ}) Anglica {\itshape th} in {\itshape the}, vel \textgreek{δ} Graecorum aetatis nostrae (1).}
% AB01-PDF0072-body-L036
\SourceLine{AB01-PDF0072-body-L036}{{\itshape ǵ} (\textarabic{\upshape ج}) semper palatalis, ut {\itshape g} Italorum ante {\itshape i} et {\itshape e}, vel Anglorum in {\itshape gentleman}.}
% AB01-PDF0072-body-L037
\SourceLine{AB01-PDF0072-body-L037}{{\itshape gh} (\textarabic{\upshape غ}) gutturalis sonora spirans, paulo asperior quam \textgreek{γ} Graecorum aetatis nostrae.}
% AB01-PDF0072-body-L038
\SourceLine{AB01-PDF0072-body-L038}{{\itshape kh} (\textarabic{\upshape خ}) paulo asperior quam {\itshape ch} Germanica in {\itshape ach}, aut {\itshape j} Hispanorum.}
% AB01-PDF0073-body-L001
\SourceLine{AB01-PDF0073-body-L001}{{\itshape sh} (\textarabic{\upshape ش}) ut apud Anglos.}
% AB01-PDF0073-body-L002
\SourceLine{AB01-PDF0073-body-L002}{{\itshape th} (\textarabic{\upshape ث}) ut Anglorum {\itshape th} in {\itshape thank}, {\itshape z} Hispanica, \textgreek{ϑ} Graecorum.}
% AB01-PDF0073-body-L003
\SourceLine{AB01-PDF0073-body-L003}{{\itshape w} (\textarabic{\upshape و}) ut apud Anglos.}
% AB01-PDF0073-body-L004
\SourceLine{AB01-PDF0073-body-L004}{{\itshape ḍ, ḥ, q, ṣ, ṭ, ẓ} emphaticae, ut vulgo dicunt, sunt simpliciorum {\itshape d, h, k, s, t, z}.}
% AB01-PDF0073-body-L005
\SourceLine{AB01-PDF0073-body-L005}{Vocales longae lineola superposita indicantur ({\itshape ā, ī, ū}).}
